\section{مقدمه و معرفی پروژه}
\label{sec:introduction}

کم‌خونی (Anemia)، به معنای کاهش سطح هموگلوبین خون به زیر آستانه‌های تعریف‌شده توسط سازمان جهانی بهداشت، یکی از شایع‌ترین مشکلات سلامت عمومی در سطح جهان است و طیف وسیعی از جمعیت را، به‌ویژه در مناطق با دسترسی محدود به خدمات درمانی، تحت تأثیر قرار می‌دهد. روش استاندارد و مرجع برای تشخیص کم‌خونی، اندازه‌گیری آزمایشگاهی غلظت هموگلوبین (Hgb) از طریق نمونه‌گیری خون است؛ روشی که هرچند دقیق است، اما تهاجمی، وابسته به تجهیزات آزمایشگاهی، و از نظر زمانی و هزینه‌ای برای غربالگری گسترده و مکرر مقرون‌به‌صرفه نیست.

از دیدگاه بالینی، رنگ‌پریدگی ملتحمه‌ی پلک (Palpebral Conjunctiva) — یعنی لایه‌ی داخلی و مخاطی پلک تحتانی چشم — از دیرباز به‌عنوان یک نشانه‌ی بصری قابل‌مشاهده و مرتبط با کاهش هموگلوبین شناخته می‌شود و در معاینات بالینی سنتی نیز مورد استفاده قرار می‌گیرد. این ارتباط، زمینه‌ساز شکل‌گیری خطوط پژوهشی متعددی شده است که به‌دنبال جایگزینی یا تکمیل روش‌های تهاجمی تشخیص کم‌خونی با روش‌های غیرتهاجمی، مبتنی بر تصویربرداری و یادگیری ماشین هستند؛ رویکردی که به دلیل هزینه‌ی پایین، سرعت بالا، و امکان استفاده در مقیاس گسترده (از جمله از طریق دوربین گوشی‌های هوشمند)، پتانسیل قابل‌توجهی برای غربالگری اولیه‌ی کم‌خونی، به‌ویژه در مناطق کم‌برخوردار، دارد.

پروژه‌ی حاضر با عنوان \lr{EYES-DEFY-ANEMIA} در همین راستا و بر پایه‌ی یک مجموعه‌داده‌ی واقعی شامل ۲۱۷ بیمار از دو کشور هند و ایتالیا شکل گرفته است. برای هر بیمار، علاوه بر مقدار هموگلوبین اندازه‌گیری‌شده (که مبنای برچسب‌گذاری دودویی کم‌خونی بر اساس آستانه‌های سازمان جهانی بهداشت است)، دو نوع تصویر ثبت شده: یک عکس خام از چشم، و یک تصویر برش‌خورده از ناحیه‌ی ملتحمه‌ی پلک. این پروژه در قالب دو فاز اصلی و مکمل تعریف شده است: فاز اول به \textbf{بخش‌بندی} (Segmentation) ناحیه‌ی ملتحمه‌ی پلک درون تصویر خام چشم می‌پردازد، و فاز دوم — که موضوع این گزارش است — به \textbf{طبقه‌بندی} (Classification) دودویی کم‌خونی/غیرکم‌خونی می‌پردازد.

\subsection{هدف فاز طبقه‌بندی}
\label{subsec:classification-goal}

هدف اصلی فاز چهارم پروژه، طراحی، پیاده‌سازی و مقایسه‌ی دقیق و منصفانه‌ی چندین معماری یادگیری عمیق برای طبقه‌بندی دودویی کم‌خونی، مستقیماً از روی تصاویر خام چشم بیماران است. این فاز به‌صورت کاملاً مستقل از فاز بخش‌بندی طراحی و پیاده‌سازی شده — یعنی هیچ‌گونه وابستگی کدی یا داده‌ای به خروجی‌های پردازش‌شده‌ی فاز اول ندارد — تا صحت و استقلال نتایج هر دو فاز به‌طور جداگانه قابل ارزیابی باشد.

فراتر از صرفِ رسیدن به بالاترین دقت ممکن، این فاز از پروژه از ابتدا با دو دغدغه‌ی روش‌شناختی اضافی طراحی شده است:

\begin{enumerate}[rightmargin=0.5cm]
    \item \textbf{مقایسه‌ی گسترده و منصفانه‌ی معماری‌ها.} به‌جای تکیه بر یک معماری واحد، مجموعه‌ای متنوع از معماری‌های شبکه‌ی عصبی کانولوشنی (CNN) و مبتنی بر ترنسفورمر (Vision Transformer) تحت یک پروتکل آموزشی یکسان و با جست‌وجوی خودکار ابرپارامتر مقایسه شده‌اند، تا نتیجه‌ی نهایی نه متکی به انتخاب تصادفی یک معماری، بلکه حاصل یک مقایسه‌ی نظام‌مند باشد.
    \item \textbf{حساسیت به سوگیری جمعیت‌شناختی.} ترکیب جمعیتی مجموعه‌داده میان دو کشور هند و ایتالیا از نظر نرخ شیوع کم‌خونی به‌شدت نامتوازن است. این عدم توازن، ریسک واقعی یادگیری میان‌بر (Shortcut Learning) را ایجاد می‌کند؛ یعنی امکان اینکه مدل به‌جای یادگیری نشانه‌های بصری واقعیِ مرتبط با کم‌خونی، صرفاً کشور بیمار را به‌عنوان یک نشانه‌ی جانبی یاد بگیرد. به همین دلیل، در سراسر این پروژه، معیارهای ارزیابی همواره به‌صورت تفکیک‌شده بر اساس کشور نیز گزارش و بررسی شده‌اند.
\end{enumerate}

\subsection{ماهیت تکرارشونده‌ی فرایند توسعه}
\label{subsec:iterative-nature}

پایپ‌لاین نهایی که در این گزارش تشریح می‌شود، حاصل یک فرایند تک‌مرحله‌ای و از پیش تعیین‌شده نیست، بلکه محصول چندین دور توسعه‌ی تکرارشونده است: آغاز با یک آزمایش اولیه بر روی سه معماری پایه، گسترش به مجموعه‌ای از نُه معماری متنوع (شش شبکه‌ی کانولوشنی و سه ترنسفورمر) بر روی دو نوع برش تصویری، بررسی عمیق‌تر پایداری نتایج از طریق اعتبارسنجی متقاطع، به‌کارگیری ابزارهای تفسیرپذیری مدل (Grad-CAM) برای بررسی مستقیم رفتار مدل، کشف و رفع یک نقص واقعی و تا پیش از این ناشناخته در مرحله‌ی پیش‌پردازش داده، و در نهایت بازآموزی کامل مدل‌ها بر روی داده‌ی اصلاح‌شده. جزئیات کامل هر یک از این مراحل، تصمیمات فنی متناظر با آن‌ها، و یافته‌های حاصل، در بخش‌های بعدی این گزارش ارائه خواهند شد.

\subsection{ساختار گزارش}
\label{subsec:report-structure}

ادامه‌ی این گزارش در چهار بخش اصلی سازمان‌دهی شده است. بخش~\ref{sec:data-pipeline} به شرح کامل پایپ‌لاین آماده‌سازی داده، شامل نحوه‌ی بارگذاری تصاویر، افزایش داده (Data Augmentation)، و به‌طور خاص تکنیک‌های به‌کاررفته برای کاهش نویز و یکسان‌سازی پس‌زمینه‌ی تصاویر می‌پردازد. بخش~\ref{sec:architectures} معماری‌های یادگیری عمیق پیاده‌سازی و آزمایش‌شده را به‌تفصیل معرفی می‌کند. بخش~\ref{sec:implementation} بر جزئیات کلیدی پیاده‌سازی و قطعه‌کدهای اصلی \lr{PyTorch} تمرکز دارد. سرانجام، بخش~\ref{sec:results} به مقایسه‌ی دقیق مدل‌ها بر اساس معیارهای ارزیابی مناسب، و ارائه‌ی نتیجه‌گیری نهایی درباره‌ی بهترین مدل و دلایل برتری آن اختصاص دارد.
